%% kapa: exam=1 round=36 year=2025 period=2 form=A subject=appraisal_law subject_ko=감정평가관계법규 seq=22 number=22
%% source: Q-Net 감정평가사 기출문제 (cst00309) artlSeq=5231525 file=raw/1008/5231525/2025년 제36회 감정평가사 1차시험 2교시 문제지 원본.pdf
%% conversion: kordoc 4.24.51 -> tools/gen1.py (auto); review=proofread-sonnet; flags=none
\begin{question}{22}
국유재산법령상 일반재산에 관한 설명으로 옳지 않은 것은?
\begin{choices}
  \item 일반재산은 ｢국토의 계획 및 이용에 관한 법률｣ 등 다른 법률에 따라 그 처분이 제한되는 경우에 매각할 수 없다.
  \item 일반재산은 국유재산관리기금의 운용계획에 따라 국유재산관리기금의 재원으로 개발할 수 있다.
  \item 용도를 지정하여 일반재산을 매각하는 경우에는 그 재산의 매각일부터 10년 이상 지정된 용도로 활용하여야 한다.
  \item 총괄청이 일반재산을 민간사업자와 공동으로 개발한 경우 일반재산의 대부기간은 20년 이내로 할 수 있으며, 20년의 범위에서 두 차례 연장할 수 있다.
  \item 소규모 일반재산을 한 곳에 모아 관리함으로써 재산의 효용성을 높이기 위하여 필요한 경우에는 일반재산인 토지와 사유재산인 토지를 교환할 수 있다.
\end{choices}
\end{question}
\answer{4}
